% Generated draft: requires mathematical and source review.
\documentclass[11pt]{article}
\usepackage[margin=25mm]{geometry}
\usepackage{fontspec}
\setmainfont{Latin Modern Roman}
\usepackage{amsmath,amssymb,mathtools,unicode-math}
\setmathfont{Latin Modern Math}
\usepackage{enumitem,needspace,xurl,hyperref}
\setlength{\parindent}{0pt}
\setlength{\parskip}{5pt}
\setlength{\emergencystretch}{3em}
\newcommand{\N}{\mathbb{N}}
\newcommand{\Z}{\mathbb{Z}}
\newcommand{\Q}{\mathbb{Q}}
\newcommand{\R}{\mathbb{R}}
\newcommand{\C}{\mathbb{C}}
\newcommand{\F}{\mathbb{F}}
\newcommand{\A}{\mathbb{A}}
\newcommand{\PP}{\mathbb{P}}
\newcommand{\E}{\mathbb{E}}
\newcommand{\eps}{\varepsilon}
\newcommand{\dd}{\,\mathrm{d}}
\begin{document}
\section*{UnsolvedMath ID: 701345}
\textbf{Source problem code:} catalogue entry
\subsection{Definitions and mathematical statement}
Let \(\pi\) be an arbitrary set of primes, possibly empty, and let \(\pi'\) be its complement in the set of all primes. A positive integer is a \(\pi\)-number if every prime dividing it belongs to \(\pi\); in particular, \(1\) is a \(\pi\)-number. A finite group is a \(\pi\)-group if its order is a \(\pi\)-number. A subgroup \(H\leq X\) of a finite group \(X\) is a Hall \(\pi\)-subgroup if \(|H|\) is a \(\pi\)-number and \([X:H]\) is a \(\pi'\)-number.

Use the standard interpretation that a finite group \(X\) is a \(D_\pi\)-group if it possesses a Hall \(\pi\)-subgroup, any two Hall \(\pi\)-subgroups of \(X\) are conjugate in \(X\), and every \(\pi\)-subgroup of \(X\) is contained in a Hall \(\pi\)-subgroup of \(X\). The supplied source refers to Archive problem 3.62 for this terminology, but does not include that entry; agreement with its precise convention requires review.

For each \(\pi\), classify, up to isomorphism, all nontrivial finite simple \(D_\pi\)-groups \(G\) satisfying each of the following conditions separately:
\begin{enumerate}
\item[(a)] Every subgroup of \(G\) is a \(D_\pi\)-group:
\[
\forall H\leq G,\qquad H\text{ is a }D_\pi\text{-group}.
\]
This part is attributed in the source to H. Wielandt.
\item[(b)] Every subgroup of \(G\) that possesses a Hall \(\pi\)-subgroup is a \(D_\pi\)-group:
\[
\forall H\leq G,\qquad
\bigl(H\text{ possesses a Hall }\pi\text{-subgroup}\bigr)
\Longrightarrow H\text{ is a }D_\pi\text{-group}.
\]
\end{enumerate}
Here the same set \(\pi\) is used for \(G\) and every subgroup under consideration. The source says finite simple groups without imposing a nonabelian hypothesis.

The supplied source reports that all finite simple \(D_\pi\)-groups are known, citing D. O. Revin (2008). Its comment of 2013 reports that the alternating and sporadic groups satisfying each of (a) and (b) were determined by N. Ch. Manzaeva (2012). Its comment of 2025 reports that the simple groups of Lie type of rank \(1\) satisfying (a) are known, citing D. O. Revin and V. D. Shepelev (2024) and V. D. Shepelev (2025). These are supplied historical qualifications, not independently verified literature assessments.
\subsection{Short English statement}
For an arbitrary fixed set of primes \(\pi\), classify the finite simple \(D_\pi\)-groups whose subgroups all satisfy the \(D_\pi\) conditions, and separately those whose subgroups satisfy these conditions whenever they possess a Hall \(\pi\)-subgroup. Preserve the source's reported partial classifications; the referenced definition in Archive 3.62 still requires checking.
\subsection{Sources}
\begin{itemize}
\item Supplied local UnsolvedMath registry record, permanent ID \(701345\), exact source problem code \texttt{catalogue entry}, title \emph{Kourovka Notebook Problem 17.43}. The record supplies the status \texttt{open}; this label is not an independently verified mathematical assessment. Record updated \texttt{2026-10-16T14:00:00Z}.
\item E. P. Vdovin and D. O. Revin, proposers of Problem 17.43, \emph{The Kourovka Notebook (Unsolved Problems in Group Theory)}, 17th issue (2010). The supplied background also identifies the 21st issue edition (2026), arXiv:1401.0300v48. The collection description names E. I. Khukhro and V. D. Mazurov as editors. Part (a) is attributed to H. Wielandt.
\item \emph{The Kourovka Notebook}, Archive problem 3.62, cited by the source for the definition of a \(D_\pi\)-group; its text is not supplied.
\item D. O. Revin, \emph{Algebra Logic}, \textbf{47}, no. 3 (2008), 210--227. Cited by the supplied source for the classification of finite simple \(D_\pi\)-groups. No article title or link is supplied.
\item N. Ch. Manzaeva, \emph{Siberian Electron. Math. Rep.}, \textbf{9} (2012), 294--305 (Russian). The supplied comment of 2013 attributes the alternating and sporadic cases of both parts to this work. No article title or link is supplied.
\item D. O. Revin and V. D. Shepelev, \emph{Siberian Math. J.}, \textbf{65}, no. 5 (2024), 1187--1194; V. D. Shepelev, \emph{Siberian Math. J.}, \textbf{66}, no. 4 (2025), 1049--1062. The supplied comment of 2025 cites these works for simple groups of Lie type of rank \(1\) satisfying (a). No article titles or links are supplied.
\end{itemize}
\end{document}
